\chapter{Secure messaging and the debit}
\label{ch:secmsg}

\chapterbrief
  {Read the balance and spend money.}
  {Growing a class as its responsibilities grow.}
  {\textbf{Money moving off a real card.}}

\section{The problem}

You are authenticated. But the card will no longer accept a bare command --- each
one must now prove three things:

\begin{enumerate}
  \item \textbf{Authenticity} --- an 8-byte CMAC computed with $K_{mac}$.
  \item \textbf{Freshness} --- a command counter, so no frame can be replayed.
  \item \textbf{Confidentiality} (for values) --- the amount encrypted with $K_{enc}$.
\end{enumerate}

That is ``CommMode.Full'', and it is what a value file configured for full
protection demands.

\begin{analogy}{The cheque}
The \textbf{amount} is written in the box --- the encrypted data. The
\textbf{signature} proves you authorised it --- the MAC. And the \textbf{cheque
number} stops the payee photocopying it and cashing it twice --- the counter.

Remove any one and the system fails: no signature, anyone writes cheques; no
number, one cheque pays forever; no amount protection, someone alters 10 to 1000.
\end{analogy}

\section{The command counter, and the ordering that will bite you}

\code{CmdCtr} is a 16-bit counter that starts at \textbf{0} immediately after
authentication and increments by one per accepted command. It goes into every MAC
and every IV, so no two commands in a session look alike.

Here is the rule that costs people an evening:

\begin{center}
\begin{tikzpicture}[font=\scriptsize\sffamily, x=1mm, y=1mm]
  \draw[brewergrey!60, -{Stealth[length=4pt]}] (0,0) -- (108,0);
  \foreach \x/\lbl in {6/{\code{ctr = 0}}, 58/{\code{ctr = 1}}}{
     \node[font=\scriptsize\sffamily\bfseries, color=brewernavy, anchor=west] at (\x,7) {\lbl};}
  % events
  \foreach \x/\lbl/\c in {
      8/{build command MAC}/brewernavy,
      30/{send}/brewergrey,
      48/{card accepts}/breweramber!80!black,
      62/{\textbf{increment}}/brewerred,
      80/{verify response MAC}/brewergreen,
      96/{decrypt response}/brewergreen}{
     \fill[\c] (\x,0) circle (1.1pt);
     \draw[\c!60] (\x,0) -- (\x,-4);
     \node[color=\c, anchor=north east, rotate=32, font=\tiny] at (\x,-4.5) {\lbl};
  }
  \draw[brewerred, dashed] (62,6) -- (62,-2);
\end{tikzpicture}
\end{center}

\begin{trap}{The increment happens \emph{before} you verify the response}
The \textbf{command} MAC uses the pre-increment value; the \textbf{response} MAC
and response IV use the post-increment value.

Get this backwards and your command MACs are perfect while every response
verification fails --- which looks exactly like a broken MAC implementation,
sending you back to debug code that is correct.
\end{trap}

\section{The four building blocks}

\subsection{Truncated MAC}

DESFire transmits only 8 of the 16 CMAC bytes: the \textbf{odd-indexed} ones,
\code{mac[1], mac[3], ... mac[15]}. It halves the frame overhead, and 64 bits is
still ample against forgery. You can verify this one by eye:

\begin{trace}
truncate(00112233445566778899aabbccddeeff) = 1133557799bbddff
\end{trace}

\subsection{MAC input layout}

Both directions have the same shape, differing only in byte 0:

\begin{trace}
command:   [cmd ][ctrLo][ctrHi][TI0..TI3][data...]
response:  [0x00][ctrLo][ctrHi][TI0..TI3][data...]
                    ^
              2 bytes, LITTLE-endian
\end{trace}

Byte 0 of the response input is \hex{0x00} --- the success status code. You only
check a response MAC after a successful status, so it is always zero.

\subsection{IV derivation}

IVs are computed, not chosen. Each is an ECB encryption of a block built from
labels, TI and the counter:

\begin{trace}
command IV  = AES-ECB(Kenc, [0xA5, 0x5A, TI0..TI3, ctrLo, ctrHi, 00 x 8])
response IV = AES-ECB(Kenc, [0x5A, 0xA5, TI0..TI3, ctrLo, ctrHi, 00 x 8])
                             \____ the labels are SWAPPED ____/
\end{trace}

Because the counter is in there, \textbf{every command gets a unique IV for
free}. That is the property CBC needs, and it falls out of the counter rather
than requiring a separate random value.

\subsection{ISO 7816-4 padding}

Append \hex{0x80}, then zeros to the next 16-byte boundary.

\begin{trap}{Padding always adds at least one byte}
A 16-byte payload becomes 32, not 16. That is what makes the padding
unambiguously removable --- if it sometimes added nothing, you could not tell a
message ending in \hex{80} from a padded one.

\begin{jcode}
int padded = (data.length | 0x0F) + 1;   // next multiple of 16, always > length
\end{jcode}
\end{trap}

\section{Build it, part 1: grow \texttt{DesfireSession}}

The session started as ``two keys''. It now becomes the full secure-messaging
engine, because that is where the state and the keys already live.

\begin{jcode}
    // ---- EV2 secure-messaging state ----
    private final byte[] ti;   // transaction id from auth, 4 bytes
    private int cmdCounter;    // starts at 0, one per command

    // Old constructor: keys only. TI defaults to zeros -- key derivation never
    // uses TI, so the Chapter 8 test keeps passing untouched.
    // (Constructor chaining: this(...) delegates to the fuller one.)
    public DesfireSession(byte[] authKey, byte[] rndA, byte[] rndB) {
        this(authKey, rndA, rndB, new byte[4]);
    }

    public DesfireSession(byte[] authKey, byte[] rndA, byte[] rndB, byte[] ti) {
        if (authKey.length != 16) throw new IllegalArgumentException("key: 16 bytes");
        if (ti.length != 4)       throw new IllegalArgumentException("TI: 4 bytes");
        this.keyEnc = deriveSessionKey(authKey, rndA, rndB, SessionKeyType.ENC);
        this.keyMac = deriveSessionKey(authKey, rndA, rndB, SessionKeyType.MAC);
        this.ti = ti.clone();
        this.cmdCounter = 0;
    }

    // ---- the MAC layer ----

    /** Tick the counter. Wraps at 16 bits. Ordering is load-bearing. */
    public void incrementCounter() {
        // TODO: cmdCounter = (cmdCounter + 1) & 0xFFFF;
    }

    /** Keep the ODD-indexed bytes: mac[1], mac[3], ... mac[15]. */
    private static byte[] truncateMac(byte[] mac16) {
        // TODO: byte[8]; out[i] = mac16[i * 2 + 1];
        return new byte[8];
    }

    /** MAC input = [cmd][ctrLo][ctrHi][TI 4][data...], then CMAC + truncate. */
    public byte[] commandMac(byte cmd, byte[] data) {
        // TODO: byte[] input = new byte[7 + data.length];
        //       input[0] = cmd;
        //       input[1] = (byte)(cmdCounter & 0xFF);
        //       input[2] = (byte)((cmdCounter >> 8) & 0xFF);
        //       copy ti into [3..7), data into [7..);
        //       return truncateMac(new AesCmac(keyMac).compute(input));
        return new byte[8];
    }

    /** Same shape, but input[0] = 0x00 (the success status code). */
    public byte[] responseMac(byte[] data) {
        // TODO
        return new byte[8];
    }

    public boolean verifyResponseMac(byte[] data, byte[] mac8) {
        // TODO: Arrays.equals(responseMac(data), mac8)
        return false;
    }

    // ---- the crypto layer ----

    public byte[] commandIv()  { /* TODO: calcIv(0xA5, 0x5A) */ return new byte[16]; }
    public byte[] responseIv() { /* TODO: calcIv(0x5A, 0xA5) */ return new byte[16]; }

    /** Only the two label bytes differ between the two IVs. */
    private byte[] calcIv(byte b0, byte b1) {
        // TODO: byte[16] block; block[0]=b0; block[1]=b1;
        //       ti into [2..6); block[6]=ctrLo; block[7]=ctrHi; rest zero;
        //       return Aes.ecbEncrypt(keyEnc, block);
        return new byte[16];
    }

    /** ISO 7816-4: append 0x80 then zeros. ALWAYS adds >= 1 byte. */
    private static byte[] padIso7816(byte[] data) {
        // TODO: int padded = (data.length | 0x0F) + 1; ...
        return data;
    }

    public byte[] encryptCommandData(byte[] plaintext) {
        // TODO: Aes.cbcEncrypt(keyEnc, commandIv(), padIso7816(plaintext))
        return new byte[0];
    }

    public byte[] decryptResponseData(byte[] ciphertext) {
        // TODO: Aes.cbcDecrypt(keyEnc, responseIv(), ciphertext)
        //   Padding is left in place -- the caller knows its payload length.
        return new byte[0];
    }
\end{jcode}

\begin{concept}{Put state and its operations in the same place}
Note what just happened to this class. It began as a value object holding two
keys. It is now a stateful engine with a counter, an identity, and a dozen
operations.

That growth is \emph{correct} --- the new responsibilities all operate on state it
already owned. Had you instead put \code{cmdCounter} in \code{DesfireCard} and
passed it into every method, the counter would live in one class and the keys in
another, and every call site would be responsible for keeping them in step. The
bug where someone forgets to increment would be \emph{possible}.

Put the state and its operations together, and whole categories of bug become
unrepresentable.
\end{concept}

\subsection{Verify the layers}

Session built from Chapter~\ref{ch:sesskeys}'s inputs, plus \code{TI = 11223344}:

\begin{center}\small
\begin{tabular}{@{}L{0.40\linewidth}cL{0.40\linewidth}@{}}
\toprule
\thead{Call} & \thead{ctr} & \thead{Expected} \\
\midrule
\code{truncateMac(0011...eeff)} & --- & \hex{1133557799bbddff} \\
\code{commandMac(0x6C, \{01\})}  & 0 & \hex{4d131ee17b482b16} \\
\code{responseMac(\{\})}         & 1 & \hex{6b082e8062d1a89d} \\
\code{commandIv()}               & 0 & \hex{dd42134cd6970980d22066cb00ab1b16} \\
\code{encryptCommandData(96000000)} & 0 & \hex{a2eb8f000bf3ddfe6e18e8859b284591} \\
\code{responseIv()}              & 1 & \hex{cb73ae774aafd81868062906135969b9} \\
\bottomrule
\end{tabular}
\end{center}

\begin{jcode}
DesfireSession s = new DesfireSession(authKey, rndA, rndB, h("11223344"));
assertEquals("4d131ee17b482b16", hex(s.commandMac((byte) 0x6C, h("01"))));
assertEquals("dd42134cd6970980d22066cb00ab1b16", hex(s.commandIv()));
assertEquals("a2eb8f000bf3ddfe6e18e8859b284591",
             hex(s.encryptCommandData(h("96000000"))));
s.incrementCounter();                                   // ctr -> 1
assertEquals("6b082e8062d1a89d", hex(s.responseMac(new byte[0])));
assertEquals("cb73ae774aafd81868062906135969b9", hex(s.responseIv()));
\end{jcode}

\hex{96 00 00 00} is 150 little-endian --- a \texteuro1.50 debit amount.

\section{Build it, part 2: the commands}

\begin{trap}{First, a one-line fix in Chapter~\ref{ch:auth}'s code}
Your \code{authenticateEv2First} currently builds the session \emph{without} the
TI:

\begin{jcode}
session = new DesfireSession(key, rndA, rndB);        // TI = 0000 !
\end{jcode}

Change it to:

\begin{jcode}
session = new DesfireSession(key, rndA, rndB, ti);
\end{jcode}

Miss this and every MAC from here on is computed against the wrong TI. The
symptom is maddening: authentication succeeds, and then the very first
\code{getValue} fails its response MAC --- one command away from the cause.
\end{trap}

\subsection{\texttt{getValue} (\hex{0x6C}) --- read the balance}

\begin{trace}
command : 90 6C 00 00 09 [fileNo][cmdMac 8] 00
response: [encrypted 16][respMac 8] 91 00
\end{trace}

\begin{jcode}
    public int getValue(byte fileNo) throws SpiException {
        // TODO 1: byte[] header = { fileNo };
        //         byte[] mac = session.commandMac((byte) 0x6C, header);  // ctr 0
        //         send header ++ mac (9 bytes) via wrapApdu(0x6C, ...)
        // TODO 2: session.incrementCounter();   // BEFORE reading the response
        // TODO 3: drop the 2 status bytes; body must be >= 24;
        //         encrypted = body[0..16), respMac = body[16..24)
        // TODO 4: if (!session.verifyResponseMac(encrypted, respMac))
        //             throw new SpiException("GetValue response MAC failed");
        //
        //   Do NOT skip this. An unverified balance is a number an attacker
        //   chose. The whole point of the MAC is refusing to trust it.
        //
        // TODO 5: byte[] plain = session.decryptResponseData(encrypted);
        //         return decodeValue(plain);    // first 4 bytes, LE signed
        return 0;
    }

    /** Signed 32-bit little-endian from the first 4 bytes. */
    private static int decodeValue(byte[] b) {
        return (b[0] & 0xFF) | ((b[1] & 0xFF) << 8)
             | ((b[2] & 0xFF) << 16) | ((b[3] & 0xFF) << 24);
    }
\end{jcode}

\subsection{\texttt{debit} (\hex{0xDC}) and \texttt{commitTransaction} (\hex{0xC7})}

\begin{trace}
debit  : 90 DC 00 00 19 [fileNo][enc amount 16][cmdMac 8] 00  ->  [respMac 8] 91 00
commit : 90 C7 00 00 08 [cmdMac 8] 00                        ->  [respMac 8] 91 00
\end{trace}

Note the asymmetry: \textbf{debit is Full} (the amount is encrypted),
\textbf{commit is MAC-only} (it carries no data).

\begin{jcode}
    /** Stage a debit. NOTHING MOVES until commitTransaction() succeeds. */
    public void debit(byte fileNo, int amountCents) throws SpiException {
        // TODO 1: byte[] header     = { fileNo };
        //         byte[] ciphertext =
        //             session.encryptCommandData(encodeValue(amountCents));
        //         byte[] macInput   = header ++ ciphertext;  // fileNo stays PLAIN
        //         byte[] mac        = session.commandMac((byte) 0xDC, macInput);
        //         send macInput ++ mac
        // TODO 2: requireOk(response); session.incrementCounter();
        // TODO 3: respMac = the last 8 bytes BEFORE the status word;
        //         verify over EMPTY data (the reply carries none).
    }

    /** Make the staged debit permanent. */
    public void commitTransaction() throws SpiException {
        // TODO: mac = session.commandMac((byte) 0xC7, new byte[0]);
        //       send it; requireOk; incrementCounter; verify respMac over empty.
    }

    private static byte[] encodeValue(int v) {
        return new byte[]{ (byte) v, (byte) (v >> 8),
                           (byte) (v >> 16), (byte) (v >> 24) };
    }

    /** Throw unless the status word is 91 00, naming the card's code. */
    private static void requireOk(byte[] response) throws SpiException {
        if (response.length < 2) throw new SpiException("response too short");
        int sw1 = response[response.length - 2] & 0xFF;
        int sw2 = response[response.length - 1] & 0xFF;
        if (sw1 != 0x91 || sw2 != 0x00) {
            throw new SpiException(String.format("card returned %02X%02X", sw1, sw2));
        }
    }
\end{jcode}

\section{Why Debit and Commit are separate}

This is the most important idea in the chapter, and it is a design lesson far
beyond NFC.

\code{Debit} does \textbf{not} change the balance. It stages a delta in the
card's transaction buffer. Only \code{CommitTransaction} makes it permanent.

Why? Because \textbf{a contactless card can leave the field at any instant},
mid-command, with no warning. Consider what a single atomic ``debit'' command
would mean: the card is debited, the reply is lost as the user pulls the card
away, the machine never learns it succeeded, and no coffee is dispensed. The
customer has paid for nothing and you cannot prove otherwise.

With staging, the failure mode is safe: no commit means the card silently
discards the delta and the balance is exactly what it was.

\begin{concept}{Make the ambiguous case the harmless one}
\textbf{The worst case becomes ``nothing happened'' instead of ``money
vanished''.}

That is why the kiosk must dispense only \emph{after} the commit is acknowledged,
and why \code{PaymentReceipt} in Chapter~\ref{ch:kiosk} is constructed only after
the commit returns.

Whenever you design an operation that can be interrupted, ask: what does the
interrupted state look like, and who pays for it? Then arrange for the answer to
be ``nobody''.
\end{concept}

\section{Verify --- the full offline replay}

Fresh session, so debit runs at ctr 0:

\begin{center}\small
\begin{tabular}{@{}L{0.32\linewidth}cL{0.48\linewidth}@{}}
\toprule
\thead{Step} & \thead{ctr} & \thead{Bytes} \\
\midrule
\code{debit(0x01, 150)} sends & 0 & \hex{01a2eb8f000bf3ddfe6e18e885}\newline\hex{9b28459127e5ea2cb8acad7d} \\
card's debit reply MAC        & 1 & \hex{6b082e8062d1a89d} \\
\code{commitTransaction()} sends & 1 & \hex{4dc9859e1a783ee0} \\
card's commit reply MAC       & 2 & \hex{febea1fcb55208ba} \\
\code{getValue} reply (balance 1000) & 1 & \hex{7950692fc2f01e72b3cee290f7122b5c}\newline $+$ \hex{a44a7f2359b2f512} \\
\bottomrule
\end{tabular}
\end{center}

\begin{jcode}
    // extra cases in the fake's transceive switch
    case (byte) 0x6C:                       // GetValue
        return getValueReply;
    case (byte) 0xDC:                       // Debit
        assertArrayEquals(
            h("01a2eb8f000bf3ddfe6e18e8859b28459127e5ea2cb8acad7d"),
            Arrays.copyOfRange(apdu, 5, apdu.length - 1), "debit bytes wrong");
        return debitReply;
    case (byte) 0xC7:                       // CommitTransaction
        assertArrayEquals(h("4dc9859e1a783ee0"),
            Arrays.copyOfRange(apdu, 5, apdu.length - 1), "commit MAC wrong");
        return commitReply;

    @Test
    public void debitAndCommit_matchesVectors() throws SpiException {
        DesfireCard card = new DesfireCard(this, this);
        card.selectApplication(new byte[]{ 0x01, 0x02, 0x03 });
        card.authenticateEv2First((byte) 0x00, key);    // ctr 0
        card.debit((byte) 0x01, 150);
        card.commitTransaction();
    }
\end{jcode}

Put debit/commit in a \textbf{separate} \code{@Test} from the getValue one. Each
needs a fresh session, because otherwise they disagree about the counter.

\begin{trap}{The green test that tests nothing}
After wiring the \hex{0x6C} case into your fake, it is entirely possible to have
a passing test that never calls \code{getValue}. The test goes green, you feel
good, and nothing was verified.

The assertion must be \code{assertEquals(1000, card.getValue((byte) 0x01))} ---
not merely a \code{case} in a switch. After writing any test, ask: \emph{what
would I have to break to make this fail?} If the answer is ``nothing'', it is not
a test.
\end{trap}

\section{Then: the real thing}

\begin{jcode}
int before = desfire.getValue((byte) 0x01);
System.out.printf("balance = %d cents%n", before);

desfire.debit((byte) 0x01, 10);          // stage 10 cents
desfire.commitTransaction();             // make it permanent
int after = desfire.getValue((byte) 0x01);
System.out.printf("%d -> %d cents%n", before, after);
\end{jcode}

\begin{verifybox}
\begin{trace}
balance = 510 cents (5.10 EUR)
debiting 10 cents...
510 -> 500 cents (5.10 -> 5.00 EUR)
\end{trace}
\end{verifybox}

\begin{center}\itshape
You have just moved real money with code you wrote from the SPI bytes up.
\end{center}

Every layer --- register encoding, FIFO handling, anticollision, block framing,
CMAC, key derivation, mutual authentication, secure messaging --- is yours.

\section{Traps}

\begin{traptable}
Auth OK, first \code{getValue} fails its MAC & The TI one-liner: session built without TI. \\
Command MACs right, response MACs wrong & Increment ordering. It happens \emph{before} verifying the response. \\
\hex{91 BE} (boundary error) & Debit would take the balance below the file's lower limit. \\
\hex{91 9D} (permission denied) & Authenticated with the wrong key for that file's access rights. \\
\hex{91 F0} (file not found) & The value file does not exist --- the card needs provisioning. \\
Balance unchanged after debit & You forgot \code{commitTransaction()}. Working as designed. \\
First command works, second fails & Not incrementing, or incrementing twice. \\
Everything fails after one bad run & Card left mid-transaction. Power-cycle the field, re-authenticate. \\
\end{traptable}

\section{Commit}

\begin{shell}
git add -A
git commit -m "M7d: EV2 secure messaging + GetValue/Debit/Commit -- money moves"
\end{shell}
